\documentclass[options]{philosophersimprint}
%\usepackage{sectsty}
%\linespread{1.2}
\usepackage{xunicode}
\usepackage{xltxtra}
\defaultfontfeatures{Mapping=tex-text}
\setmainfont[
Ligatures={Common}, Numbers={OldStyle}, Variant=01,
BoldFont=LinLibertine_RB.otf,
ItalicFont=LinLibertine_RI.otf,
BoldItalicFont=LinLibertine_RBI.otf
]{LinLibertine_R.otf}
\setmonofont[Scale=0.8]{DejaVuSansMono.ttf}
\usepackage{stmaryrd}
\usepackage[toc,page]{appendix}
\usepackage{amsmath,amsthm,amssymb}
\usepackage[colorlinks=true, citecolor=blue, linkcolor=blue, hyperfootnotes=false]{hyperref}
\usepackage{nameref}
\usepackage{soul}
\usepackage{frame}
\usepackage[style=apa, natbib=true]{biblatex}
\usepackage{scrextend}
\usepackage{mathrsfs}
%\addtokomafont{labelinglabel}{\sffamily}
\makeatletter
\newcommand\labelitem[2]{%
  \item[#1]\def\@currentlabel{#1}\label{#2}}
\makeatother
\usepackage[inline]{enumitem}
\usepackage{framed, csquotes}
\usepackage{tikz}
\usetikzlibrary{cd}
\tikzset{%
    symbol/.style={%
        draw=none,
        every to/.append style={%
            edge node={node [sloped, allow upside down, auto=false]{$#1$}}}
    }
}
%\usepackage{anyfontsize}
\usepackage[nodayofweek,level]{datetime}
\usepackage{etoolbox}
\patchcmd{\formatdate}{,}{}{}{}
\usepackage{enumitem, hyperref}
\makeatletter
\def\namedlabel#1#2{\begingroup
    #2%
    \def\@currentlabel{#2}%
    \phantomsection\label{#1}\endgroup
}
\makeatother
%\usepackage{anyfontsize}
\usepackage[margin=1in]{geometry}
\usepackage{pifont}% http://ctan.org/pkg/pifont
\newcommand{\cmark}{\ding{51}}%
\newcommand{\xmark}{\ding{55}}
\usepackage{forest}
\forestset{
  L1/.style={draw=black,},
  L2/.style={,edge={,line width=0.8pt}},
}
\usepackage[acronym]{glossaries}
\usepackage{stackengine}
\stackMath
\usepackage{scalerel}
\usepackage{titlesec}
\newcommand\halo{{{\circ}\mkern-1mu}}
\usepackage{linguex}
\usepackage{parskip}
\usepackage{fancyhdr}
\pagestyle{plain}
% Diagrams
%\usepackage{pgfplots}
\usepackage{tikz}
\usetikzlibrary{arrows,calc,patterns,positioning,shapes}
\usetikzlibrary{decorations.pathmorphing}
\tikzset{
modal/.style={>=stealth’,shorten >=1pt,shorten <=1pt,auto,
node distance=1.5cm,semithick},
world/.style={circle,draw,minimum size=1cm,fill=gray!15},
point/.style={circle,draw,fill=black,inner sep=0.5mm},
reflexive/.style={->,in=120,out=60,loop,looseness=#1},
reflexive/.default={5},
reflexive point/.style={->,in=135,out=45,loop,looseness=#1},
reflexive point/.default={25},
}
\usepackage{tree-dvips}
\usepackage{qtree}

\DeclareMathOperator*{\argmax}{arg\,max}
\DeclareMathOperator*{\argmin}{arg\,min}
\newtheorem{theorem}{Theorem}
\newtheorem{corollary}{Corollary}
\newtheorem{lemma}{Lemma}
\newtheorem{proposition}{Proposition}
\theoremstyle{definition}
\newtheorem{definition}{Definition}
\newtheorem{example}{Example}
\newtheorem{fact}{Fact}
\newtheorem{remark}{Remark}
\newtheorem{claim}{Claim}
\newcommand\numberthis{\addtocounter{equation}{1}\tag{\theequation}}

\newcommand{\ev}[2][]{\ensuremath{\llbracket #2\rrbracket^{#1}}}
%command
\newcommand{\A}{\mathscr{A}}
\newcommand{\cred}{\mathcal{C}}
\newcommand{\des}{\mathcal{D}}
\newcommand{\scred}{\mathsf{C}}
\newcommand{\sdes}{\mathsf{D}}
\newcommand{\ut}{\mathcal{U}}
\newcommand{\rb}{{\rrbracket}}
\newcommand{\lb}{{\llbracket}}
\newcommand{\counterfactual}{\ensuremath{%
  \Box\kern-1.5pt
  \raise1pt\hbox{$\mathord{\rightarrow}$}}}
  \newcommand{\tr}{\mathsf{T}}
%\renewcommand{\fa}{\mathsf{F}}

\usepackage{FOL_Yale/notation}

\begin{document}

PHIL 1115: First-order Logic \hfill Yale, Fall 2026\\
Lecture 1, Handout \hfill Dr. Daniel Grimmer

\textbf{1. What is Logic? What is ``Good Reasoning''?}

\textbf{Discussion Question}. If asked on the street, how would you define logic?

At least in part, \textbf{logic} is the study of arguments.\footnote{Notation: in these notes, I'll use \textbf{bold-face} text when I'm defining, or at least characterizing, a technical term for the first time.} In particular, logic aims to tell us which are the ``good'' arguments and which are the ``bad'' arguments.

\textbf{Discussion Question}. If asked on the street, how would you characterize a ``good argument'' or ``good reasoning''? What is ``bad reasoning''?

Throughout this course, we'll think of an \textbf{argument} as consisting of some number of \textbf{premises} (maybe many, maybe none) together with a \textbf{conclusion}. Usually, we signal the conclusion with words like `So' or `Therefore'. It is marked with the symbol `$\therefore$' read as `therefore'. Here is an example:

\textbf{Argument 1 (Good, Disjunctive Syllogism).}
\begin{itemize}
    \item[] Either the butler did it or the gardener did it.
    \item[] The butler did not do it.
    \item[$\therefore$] The gardener did it.
\end{itemize}
Intuitively, this is a good argument. If we know that the premises are both true (e.g., by our rigorous detective work), then the conclusion must also be true.

But suppose that there is a twist in the case: the gardener is innocent, they have a solid alibi! Our conclusion is wrong.

\textbf{Discussion Question}. Is this still a good argument? Can good reasoning lead us to a false conclusion?

Arguably, yes. It remains good because \textit{if the argument's premises were both true, then the conclusion would have also been true.} Said differently, if a good argument's conclusion turns out to be wrong, then there must be an issue with one of its premises; namely, one of its initial premises must have been wrong. In our case, this means that: either the butler did it, or we were wrong to limit ourselves to only these two suspects.

\textbf{Key Thought.} There are two stages in which an argument can go wrong: bad premises or bad reasoning (i.e., a bad premise-to-conclusion linkage).

Here is an argument with good premises but bad reasoning.

\textbf{Argument 2 (Bad, Affirming the Consequent).}
\begin{itemize}
    \item[] If it is raining, then the grass is wet.
    \item[] The grass is wet.
    \item[$\therefore$] It is raining.
\end{itemize}
Could it ever be that the premises are both true and yet the conclusion is false? Sure, the grass could be wet for some other reason (sprinklers, dew, water balloon fight). This is a bad argument: The truth of the conclusion does not automatically follow from the truth of the premises.\footnote{Equivalently, the falsity of the conclusion does not automatically reveal a falsity among the premises.}

What desiderata (D) do we expect from a perfectly good argument? Intuitively,
\begin{enumerate}
    \item[D1:] it should start from true premises.
    \item[D2:] it should produce a true conclusion.
    \item[D3:] the truth of the premises should somehow \textit{guarantee} the truth of the conclusion \textit{as a matter of logical necessity}.
\end{enumerate}
To see why we need this third condition consider Argument 2 in the case where it does happen to be raining. In this case we would have true premises and a true conclusion but no real ``therefore''. It is still bad reasoning even if whoever made the argument happened to get lucky.

Going forward, let's say that an argument is \textbf{sound} just in case it satisfies all three of these desiderata (D1, D2, and D3).\footnote{It is more common to define soundness in terms of validity as follows: An argument is sound just in case it is valid (D3) and has true premises (D1). Note that given D1 and D3, we get D2 for free. It is listed here only to boost your intuition. A warning about this word, for later: in Lecture 7 we will meet a second, unrelated notion of \textit{soundness}, applied not to an argument but to a whole method of proof. Keep the two apart.} A sound argument is one with true premises and a true conclusion with the premises somehow guaranteeing the truth of the conclusion. Sound arguments are the gold standard.

Interestingly, however, soundness is not the central concept in logic. Far more important is the notion of validity (just D3). An argument is \textbf{valid} just in case it's logically impossible for its premises to all be true and yet for its conclusion to be false. Intuitively, valid arguments are ones whose conclusion automatically ``follows from'' the premises as a matter of logical necessity. Here is an epistemic gloss: anyone who accepts the premises \textit{must} accept the conclusion, e.g., on pain of ``being illogical''. (We'll refine these definitions later on.)

But is it really true that logicians only care about validity rather than soundness? Yes: \textit{Logicians do not care about which particular claims happen to be true or false!}\footnote{Arguably, there is one kind of truth which logicians do care about: tautologies (a.k.a. theorems, logical truths). More on these in Lecture 4.} For instance, checking the soundness of Argument 2 would involve the logician getting out of their armchair to check whether or not the grass is wet. Checking the truth of individual claims is not a part of the logician's job description.

Similarly, consider an argument in the context of either detective work (Argument 1) or biology. In order to check that such an argument is sound, one must check that its premises are true. But establishing such truths is the job of the detective or the biologist, not the logician. Similarly for every subject. Arguably, even mathematics! (This point is much debated.)

Instead, the logician's subject matter is \textit{validity}. What constitutes a good premise-to-conclusion linkage? Something about ``good reasoning'' makes it the case that:
\begin{enumerate}
    \item whenever the premises are true, \textit{the flawless structure of the argument somehow guarantees} that the conclusion is also true.
    \item Equivalently, whenever the conclusion is false \textit{the flawless structure of the argument somehow guarantees} that we can trace its falseness to at least one premise being false.
\end{enumerate}

\textbf{Key Thought.} Validity is whatever feature of a good argument allows for it to faithfully carry truth from its premises to its conclusion. (And equivalently, to carry the falsity of its conclusion back to one of its premises).

\textbf{2. Formal Validity - Truth Preserving by Virtue of its Form}

Let's see an example of an argument which is valid but not sound.

\textbf{Argument 3.}
\begin{itemize}
    \item[] All lizards are mammals.
    \item[] No mammals are egg-layers.
    \item[$\therefore$] No lizards are egg-layers.
\end{itemize}
Any biologist will tell you that all three of these sentences are false. Nonetheless, the logician considers this to be a good argument because of its validity: \textit{If the premises would have been true, then the conclusion would have been true as well.}

But how can we establish the validity of this argument? Doesn't it require thinking through the biological details of a bizarre counterfactual world in which these premises are true? Our logician would hate that; it is not their job to do biology! It is not their job to do anything except logic. In fact, our logician has never even seen the words ``lizards'', ``mammals'', and ``egg-layers'', and has no idea what they mean. They just look like squiggles. So let's replace them with other squiggles.

\textbf{Form of Argument 3 (Celarent, EAE-1. Aristotle's Prior Analytics)}
\begin{itemize}
    \item[] All Xs are Ys.
    \item[] No Ys are Zs.
    \item[$\therefore$] No Xs are Zs.
\end{itemize}
It should now be clear that establishing the validity of this argument does not require any biological knowledge whatsoever.

\textbf{Key Thought}. Some arguments are valid purely in virtue of their \textit{form}.\footnote{A caution you will also meet in this week's reading: an argument has \textit{many} forms, not one. Argument 3 equally instantiates ``First thing; Second thing; $\therefore$ Third thing'', which has invalid instances aplenty. So ``valid in virtue of its form'' means that \textit{some} form of it has only valid instances, and finding that form is what this course is for.}

This is what is meant by ``formal'' logic. The rest of this course will be spent developing this idea: An argument is good (i.e., valid) if it is truth preserving by virtue of its form.

\textbf{3. Some More Examples of Valid Argument Forms}

\begin{itemize}
    \item[Arg. 4] If Socrates is a man, then Socrates is mortal.
    \item[] Socrates is a man.
    \item[$\therefore$] Socrates is mortal.
\end{itemize}
\begin{itemize}
    \item[Arg. 5] If the moon is made of green cheese, then my cat is the President.
    \item[] The moon is made of green cheese.
    \item[$\therefore$] My cat is the President.
\end{itemize}
\begin{itemize}
    \item[Arg. 6] If squarks are funky, then glunks are squeem.
    \item[] Squarks are funky.
    \item[$\therefore$] Glunks are squeem.
\end{itemize}

\textbf{Discussion Question}. Which of the above three arguments are good (valid)?

Notice that all three arguments have the same form: \textit{if a first thing, then a second thing; the first thing; therefore the second thing.} Logicians have a special name for this pattern of argumentation, \textbf{modus ponens}.

\textbf{Form of Arguments 4, 5, and 6. (Modus Ponens)}
$$\text{If $p$, then $q$,\qquad $p$\qquad $\therefore$ $q$.}$$
Arguments 4, 5, and 6 are all valid: truth preserving by virtue of their form. The argument form remains valid \textit{no matter what sub-sentences we substitute for $p$ and $q$}.

What options are available when someone presents you with a modus ponens argument? A key philosophical aphorism is ``\textit{one man's modus ponens is another man's modus tollens}''. Here is a historically significant example.

\textbf{G.E. Moore vs. The Skeptic} - The skeptic argues by modus ponens:\\
\textbf{Argument 7.}
\begin{itemize}
    \item[] If skepticism is true, you cannot know that you have hands. (If $p$ then $q$)
    \item[] Skepticism is true. ($p$)
    \item[$\therefore$] You cannot know you have hands. ($q$)
\end{itemize}

Now, Moore agrees with the Skeptic about the truth of the first premise (if $p$ then $q$), but he disagrees about the conclusion ($q$). In particular, Moore argued that his common-sense certainty of having hands was much stronger than any abstract philosophical argument for skepticism.

But modus ponens is a valid argument form: Hence, if you disagree with the conclusion you must disagree with one of the premises. The only premise left is the second one, so Moore is forced to deny skepticism. Thus, in response to the skeptic's argument Moore reasons as follows.

\textbf{Argument 8.}
\begin{itemize}
    \item[] If skepticism is true, you cannot know that you have hands. (If $p$ then $q$)
    \item[] You can know you have hands: ``Here is one hand, here is another''. (Not $q$)
    \item[$\therefore$] Skepticism is false. (Not $p$)
\end{itemize}
The general form of Moore's counter argument is called \textbf{modus tollens}:

\textbf{Form of Argument 8 (Modus Tollens)}
$$\text{If $p$, then $q$,\qquad not $q$\qquad $\therefore$ not $p$.}$$
This pattern of argument and counter-argument recurs across all philosophical disciplines. Both sides of the debate accept that ``if $p$ then $q$''. One side believes $p$ (and hence $q$) whereas the other side believes ``not $q$'' (and hence must believe ``not $p$''). Again: ``\textit{one man's modus ponens is another man's modus tollens}''.

This aphorism perfectly captures the limits of formal logic. Logic can tell you if an argument is structured correctly, but it cannot generally tell you which starting premises to accept. When a modus ponens argument leads to a conclusion that one person finds entirely absurd (a ``bullet'' they refuse to bite), they will simply use modus tollens to reverse-engineer a rejection of one of the original premises. Both of these maneuvers are logically valid.

\textbf{4. Propositions and Sentences} Our characterization of valid arguments (truth-preserving by virtue of its form) requires the sentences inside the argument to be \textit{truth-apt}. They need to be capable of being true or false. In philosophy jargon, they need to be \textbf{propositions}. Philosophers disagree about what kind of thing a proposition is. To help clarify, here are some non-examples:
\begin{itemize}
    \item Was Albert Einstein born in Vienna, Austria?
    \item Take out the trash!
    \item Colorless green ideas sleep furiously.
\end{itemize}
The first two sentences here (the question and the command) don't express propositions because they don't have the right grammatical \textit{form}; while they are meaningful, they're not capable of saying something that's true or false. The last sentence, in contrast, does have the right grammatical form (it looks like a declarative sentence). Notice, however, that it still doesn't seem to express a proposition. (Indeed, it doesn't seem to \textit{express} anything at all.)

As a rule of thumb then, propositions are expressed by meaningful declarative sentences. There are, however, a few complications. In particular, some propositions can be expressed by more than one sentence:
\begin{itemize}
        \item[$\bullet$] It's raining. \qquad\qquad\qquad $\bullet$ It's precipitating.
        \item[$\bullet$] Il pleut. (French) \qquad\qquad $\bullet$ Es regnet. (German)
\end{itemize}

Conversely, the same sentence can be used to express different propositions, because it's \textit{context-sensitive}, \textit{ambiguous}, or contains \textit{indexicals}:
    \begin{itemize}
        \item Steph Curry is tall. (Context-sensitive)
        \item Later, I'll go to the bank. (Ambiguous)
        \item I'm here now. (Contains indexicals. Three, in fact)
    \end{itemize}

In this course, we'll restrict ourselves to the parts of English that are \textit{never} context-sensitive, ambiguous, etc. In the next lectures, we will develop a precise, \textit{formal} language, which ensures that ambiguity, etc., are nowhere to be found. So we don't have to get hung up on the complexities of English.

\textbf{Many Forks in the Road}. The system of logic we are about to build reflects a series of choices. For nearly every principle we adopt, there is a coherent alternative that some logician has explored, and walking through those doors often turns out to be fruitful rather than forbidden. Throughout this course I will try to flag these \textit{forks in the road} as we reach them. Logic is not a fixed list of commandments to be memorized; it is a living, contested field, and you are invited to ask, at every step: ``why \textit{this}, and not otherwise?''

For instance, our characterization of validity required that the sentences under consideration be capable of being true or false. Moreover, it requires that in each logically possible world (formally, in every model) each proposition is \textit{exactly one} of true or false (never both, never neither, never some secret third thing). To some, this may seem metaphysically heavy.

In Lecture 9 we will introduce an alternative notion of validity based on logical proofs (the natural deduction system). There, an argument is considered valid if one can use logical step-by-step reasoning to get from its premises to its conclusion. This alternative conception of validity does not require that we think in terms of truth/falsity across logically possible worlds, but rather whether something has been proven (or remains unproven) in the actual world.
\end{document}
